\section{Considerações Finais}

Este artigo propõe investigar se o estoque de vegetação previamente existente modifica, em magnitude economicamente relevante, a resposta da produtividade agrícola aos choques climáticos. Ao concentrar-se em H2 e em sua extensão territorial, H2-a, o trabalho organiza uma parte específica da relação de interdependência entre economia e meio ambiente: a possibilidade de o estado ambiental alterar as condições sob as quais a produção responde a perturbações externas.

A revisão oferece fundamentos para essa investigação, mas também evidencia diferenças de escala, cultura e mecanismo que impedem pressupor um efeito uniforme da vegetação. O desenho proposto procura enfrentar essa dificuldade ao combinar rendimento de soja, cobertura anterior à janela agrícola e choques de precipitação em unidades territoriais comparáveis de Mato Grosso. A contribuição pretendida recai sobre a mensuração da resposta produtiva, de sua magnitude e de sua heterogeneidade, mantendo o sinal das interações aberto à evidência.

A execução depende da harmonização territorial, da definição do calendário agrícola e da referência climática, da revisão das classes de vegetação e do tratamento dos símbolos de produção. Persistem também limites de identificação: anterioridade não equivale a exogeneidade, chuva local pode responder à cobertura e características produtivas não observadas podem modificar a sensibilidade ao clima. O reduzido suporte do Pantanal poderá restringir a comparação econométrica entre os três biomas.

Sendo assim, não se antecipa a confirmação das hipóteses. Espera-se que a análise contribua para avaliar em que medida as condições ambientais integram os resultados produtivos e se diferenças locais devem receber maior atenção na discussão da conservação. Essa contribuição será complementar ao exame dos incentivos econômicos e institucionais realizado no capítulo 1, sem substituir a avaliação específica de políticas ou a modelagem intertemporal das escolhas dos agentes.
